%% kapa: exam=1 round=25 year=2014 period=2 form=A subject=appraisal_law subject_ko=부동산관계법규 seq=7 number=7
%% source: Q-Net 감정평가사 기출문제 (cst00309) artlSeq=5101207 file=extracted/1008/5101207/2014년 제25회 감정평가사 1차시험 기출문제/2교시(A형).hwp
%% conversion: kordoc 4.24.51 -> tools/gen1.py (auto); review=proofread-sonnet; flags=none
\begin{question}{7}
국토의 계획 및 이용에 관한 법령상 도시지역의 지구단위계획구역에서 건축물을 건축하려는 자가 1,000제곱미터의 대지 중 400제곱미터를 공공시설 부지로 제공하는 경우 그 건축물에 적용되는 최대 용적률은? (단, 당해 대지 및 공공시설 제공 부지에 적용되는 용적률은 100퍼센트이고, 용적률의 상한은 고려하지 않음)
\begin{choices}
  \item 140퍼센트
  \item 160퍼센트
  \item 180퍼센트
  \item 200퍼센트
  \item 220퍼센트
\end{choices}
\end{question}
\answer{4}
